% Generated draft: requires mathematical and source review.
\documentclass[11pt]{article}
\usepackage[margin=25mm]{geometry}
\usepackage{fontspec}
\setmainfont{Latin Modern Roman}
\usepackage{amsmath,amssymb,mathtools,unicode-math}
\setmathfont{Latin Modern Math}
\usepackage{enumitem,needspace,xurl,hyperref}
\setlength{\parindent}{0pt}
\setlength{\parskip}{5pt}
\setlength{\emergencystretch}{3em}
\newcommand{\N}{\mathbb{N}}
\newcommand{\Z}{\mathbb{Z}}
\newcommand{\Q}{\mathbb{Q}}
\newcommand{\R}{\mathbb{R}}
\newcommand{\C}{\mathbb{C}}
\newcommand{\F}{\mathbb{F}}
\newcommand{\A}{\mathbb{A}}
\newcommand{\PP}{\mathbb{P}}
\newcommand{\E}{\mathbb{E}}
\newcommand{\eps}{\varepsilon}
\newcommand{\dd}{\,\mathrm{d}}
\begin{document}
\section*{UnsolvedMath ID: 5000007}
\textbf{Source problem code:} catalogue entry
\subsection{Definitions and mathematical statement}
Let $D$ be the surface of a regular dodecahedron, and let $G_D$ be its edge graph. For vertices $v,w$, the graph distance $d_{G_D}(v,w)$ is the minimum number of edges in an edge path joining them; it is not intrinsic surface distance.

A geodesic segment $\gamma:[0,L]\to D$ beginning at $v$ is locally straight in the planar unfolding of successive faces. In the source's terminology it is \emph{short} if its endpoint $w=\gamma(L)$ is a vertex and no interior point $\gamma(t)$, $0<t<L$, is a vertex. This is an interior-vertex exclusion, not a numerical length bound or an assertion that the segment globally minimizes distance.

The source calls such a segment type $A_0$ when its endpoint-angle parameters satisfy
\[
 \beta-\alpha=\frac{2\pi}{5}.
\]
However, it does not define $\alpha$ and $\beta$, their reference edges/rays, orientations, or branch choices in the supplied excerpt. These geometric conventions are needed to identify the type-$A_0$ class; no arbitrary angles are introduced to make the relation meaningful.

\textbf{Conjecture 3.2 (Dmitry Fuchs).} For every starting vertex $v$ and every short geodesic of type $A_0$ beginning at $v$, its terminal vertex $w$ satisfies
\[
 d_{G_D}(v,w)\ne2.
\]
The source excludes exactly graph-distance-two endpoints; it does not require the endpoint to equal $v$ or exclude other graph distances. No simplicity, absence of self-intersection, or extra length restriction is added. Allowed edge-following or degenerate endpoint directions also require the original angle/geodesic conventions.

The supplied summary says this exact endpoint exclusion was presented as open in the 2020 paper and no later resolution was found through its 2026 search. It relates the work to Fuchs--Tabachnikov billiard/geodesic studies, without giving a theorem matching this conjecture. These are supplied evidence limitations, not an independent current-status conclusion.
\subsection{Short English statement}
A short geodesic on the regular dodecahedron starts and ends at vertices and meets no other vertex. Fuchs conjectures that those of type $A_0$, specified by $\beta-\alpha=2\pi/5$, never end two edges away from their starting vertex. The angle conventions defining the type are missing.
\subsection{Sources}
\begin{itemize}
\item Local UnsolvedMath registry, permanent ID 5000007, source code catalogue entry, ``Short geodesics on the regular dodecahedron''; proposer: Dmitry Fuchs, proposed year 2020. Status open, classification OPEN-TRIAGE, updated 2026-08-06.
\item Fuchs, \emph{Billiard Trajectories in Regular Polygons and Foliations by Closed Geodesics on Surfaces} (2020), Conjecture 3.2, \url{https://amj.math.stonybrook.edu/pdf-Springer-final/020-0170.pdf}. Reported PDF-text extraction and access: 2026-07-29. The record says the source presents an open item and no later authoritative resolution check is recorded. Rights: NEEDS\_REVIEW; public access verified but redistribution permission not established.
\item The summary mentions Fuchs--Tabachnikov work and ``More on periodic billiard trajectories in the cube'' as related context, but supplies no precise result or full citation addressing this dodecahedral endpoint claim.
\end{itemize}
\end{document}
